\subsection{Обзор аналогов}
\label{sec:analysis:analogues}

Задачу отслеживания выгодных предложений уже решают несколько типов программных средств, но каждое из них рассчитано на свои условия. Для обзора выбраны три аналога, которые подходят к задаче с разных сторон: встроенный механизм самой площадки \textit{Kufar.by}, зарубежные сервисы отслеживания цен на новые товары и открытые \textit{Telegram}-боты для мониторинга объявлений. Каждый аналог рассмотрен по одной схеме: назначение, принцип работы, достоинства и недостатки применительно к рынку подержанных товаров.

\subsubsection{Сохранённые поиски \textit{Kufar.by}}

Сохранённый поиск~-- встроенная функция площадки \textit{Kufar.by} для зарегистрированных пользователей. Пользователь выбирает на сайте категорию, задаёт фильтры (бренд, характеристики товара, регион, диапазон цены) и сохраняет получившийся поиск в личном кабинете. После этого площадка оповещает владельца поиска о новых объявлениях, которые попадают под заданные условия [1].

Технически сохранённый поиск сводится к набору параметров обычной поисковой выдачи. При открытии сохранённого поиска площадка переходит по ссылке, в которой закодированы все выбранные фильтры: код категории, коды значений характеристик и границы цены. Эти же коды площадка использует во внутреннем поисковом интерфейсе. Ссылка полностью воспроизводима, поэтому набор фильтров, собранный на сайте, можно перенести в другую систему, если та понимает тот же формат параметров.

Достоинства сохранённых поисков:
\begin{itemize}
	\item функция встроена в площадку, бесплатна и не требует сторонних сервисов;
	\item доступен полный набор фильтров площадки по всем категориям;
	\item сведения о новых объявлениях поступают из первоисточника.
\end{itemize}

Недостатки проявляются, как только задача выходит за рамки «показать новые объявления»:
\begin{itemize}
	\item цена работает только как фильтр выдачи, а не как порог выгодности, относительно которого оценивается предложение;
	\item система реагирует лишь на появление новых объявлений, снижение цены на уже опубликованное объявление не отслеживается;
	\item история цен не сохраняется, и пользователь не видит, как менялась цена объявления;
	\item оповещения приходят только через сервисы площадки;
	\item программный доступ к сохранённым поискам закрыт: запрос к ним без авторизованной сессии владельца завершается ошибкой 401, и прочитать их другой программой нельзя.
\end{itemize}

\subsubsection{Сервисы отслеживания цен \textit{CamelCamelCamel} и \textit{Keepa}}

\textit{CamelCamelCamel} и \textit{Keepa}~-- наиболее известные сервисы отслеживания цен на товары интернет-магазина \textit{Amazon}. Пользователь указывает товар ссылкой на страницу товара или артикулом, после чего сервис регулярно фиксирует цену и строит график изменения цены за всё время наблюдения. Пользователь задаёт желаемую цену и получает уведомление, когда фактическая цена опускается до неё. \textit{Keepa}, кроме веб-сайта, распространяется как расширение для браузера и встраивает график истории цен прямо в страницу товара. Часть возможностей \textit{Keepa} доступна только по платной подписке [2, 3].

По модели данных эти сервисы ближе всего к разрабатываемой системе. Оба сервиса хранят историю цен, сравнивают текущую цену с порогом пользователя и уведомляют о снижении цены, то есть решают задачу управления данными о ценах, а не только показа этих данных.

Достоинства сервисов:
\begin{itemize}
	\item полная история цены с момента начала наблюдения и наглядные графики;
	\item пороговая цена как явное условие уведомления;
	\item уведомления о снижении цены до заданного уровня.
\end{itemize}

Недостатки с точки зрения рынка подержанных товаров:
\begin{enumerate}
	\item Поддерживается только \textit{Amazon}, белорусские площадки объявлений не охвачены.
	\item Модель данных опирается на артикул: один товар имеет одну цену, меняющуюся во времени. На площадке объявлений каждое предложение уникально, артикула у него нет, и история цены по артикулу к нему неприменима.
	\item Сервисы не ищут новые предложения по критериям, а следят только за заранее указанным товаром.
	\item Характеристики, важные для подержанного товара (состояние, комплектация, район), не учитываются.
\end{enumerate}

\subsubsection{Открытые \textit{Telegram}-боты мониторинга объявлений}

На \textit{GitHub} опубликованы небольшие проекты с открытым исходным кодом, которые отслеживают объявления на \textit{Kufar.by}, \textit{Avito} и \textit{OLX} и присылают их в \textit{Telegram}. Типичная реализация устроена просто. Пользователь передаёт боту поисковую фразу или ссылку на выдачу. Бот с заданным интервалом запрашивает выдачу площадки, сравнивает идентификаторы полученных объявлений со списком уже виденных и о новых отправляет сообщение. Список виденных объявлений обычно хранится в файле или в памяти процесса.

Достоинства открытых ботов:
\begin{itemize}
	\item уведомления приходят в \textit{Telegram}, которым пользователь и так пользуется каждый день;
	\item боты бесплатны, а исходный код доступен для изучения и доработки;
	\item развернуть такого бота можно за несколько минут.
\end{itemize}

Недостатки:
\begin{enumerate}
	\item Поиск ведётся по тексту, характеристики категорий не учитываются. Объявление с названием бренда, написанным кириллицей, по латинскому названию не находится, а по слову «велосипед» в выдачу попадают камеры и втулки.
	\item Сравниваются только идентификаторы объявлений, поэтому снижение цены на уже виденное объявление остаётся незамеченным.
	\item Частота запросов к площадке не ограничивается, и на отказ площадки (ответы 403 и 429) бот не реагирует, продолжая запросы в прежнем темпе до блокировки.
	\item Сообщение отправляется один раз. При отказе \textit{Telegram} уведомление теряется.
	\item Хранение в файле рассчитано на одного пользователя и не выдерживает одновременной работы нескольких процессов.
\end{enumerate}

\subsubsection{Сравнение аналогов}

Результаты обзора сведены в таблицу~\ref{table:analysis:analogues}. В последнем столбце указаны возможности разрабатываемой системы.

\begin{table}[ht]
\caption{Сравнение аналогов и разрабатываемой системы}
\label{table:analysis:analogues}
\centering
	\begin{tabular}{|>{\raggedright}m{0.20\textwidth}
	                |>{\centering}m{0.17\textwidth}
	                |>{\centering}m{0.17\textwidth}
	                |>{\centering}m{0.15\textwidth}
	                |>{\centering\arraybackslash}m{0.17\textwidth}|}
	\hline
	\centering Критерий & Сохранён\-ные поиски \textit{Kufar.by} & \textit{CamelCamel\-Camel}, \textit{Keepa} & Открытые \textit{Telegram}-боты & Разрабаты\-ваемая система \\
	\hline Работа с \textit{Kufar.by} & $+$ & $-$ & $+$ & $+$ \\
	\hline Пороговая цена как условие уведомления & частично & $+$ & частично & $+$ \\
	\hline Уведомление о снижении цены & $-$ & $+$ & $-$ & $+$ \\
	\hline Хранение цен по объявлению & $-$ & $+$ & $-$ & частично \\
	\hline Отбор по характеристикам категории & $+$ & $-$ & $-$ & $+$ \\
	\hline Уведомления в \textit{Telegram} & $-$ & $-$ & $+$ & $+$ \\
	\hline Экспорт данных и отчёты & $-$ & частично & $-$ & $+$ \\
	\hline
	\end{tabular}
\end{table}

Ни один из аналогов не закрывает задачу целиком. Сохранённые поиски площадки работают с полным набором её фильтров, но не отслеживают цену объявления во времени и закрыты для других программ. \textit{CamelCamelCamel} и \textit{Keepa} хранят историю цен и работают с порогом, но рассчитаны на новые товары с артикулом и на \textit{Amazon}. Открытые боты доставляют уведомления в \textit{Telegram}, однако ограничиваются текстовым поиском по идентификаторам и теряют сообщения при сбоях. Свободной остаётся ниша системы, которая объединяет фильтры площадки, пороговую цену, отслеживание снижения цены и надёжную доставку уведомлений в \textit{Telegram} для нескольких пользователей одновременно. Все перечисленные функции опираются на хранимые данные, и их реализация требует базы данных, а не файла со списком виденных объявлений.
